\documentclass[11pt]{article}
\usepackage{amsmath,amssymb,amsthm,mathtools,hyperref}

\begin{document}
\section*{Tumour vaccinology}
\label{prob:258}
\label{rank:42}
\par\smallskip

\noindent\textbf{Domain.} Cancer immunology; antigen selection. \quad \textbf{Retained status/source context.} Open generalization and efficacy question. Therapeutic neoantigen vaccines already elicit responses in selected trials; their existence is not an unsolved conjecture.

\paragraph{Definitions and assumptions.} For a specified patient population let $\mathcal A$ be candidate antigens, $K$ a vaccine size limit, $p_j$ the prevalence of malignant clone $j$, and $r_{ijh}$ the probability that antigen $i$ yields effective recognition of clone $j$ in HLA stratum $h$. Let $q_i$ bound normal-tissue cross-reactivity from a defined assay panel. Antigen presentation, T-cell recognition and cytotoxicity are separate measured steps. All probabilities refer to a fixed platform and immune-adjuvant context; a predicted binding score is not itself a recognition probability.

\paragraph{Problem statement.} Infer calibrated $r_{ijh}$ and select $S\subset\mathcal A$ with $|S|\leq K$ to maximize robust coverage:
\[
\min_{h\in\mathcal H}\sum_jp_j\left(1-\prod_{i\in S}(1-r_{ijh})\right),
\qquad \sum_{i\in S}q_i\leq q_{\max}.
\]
The product assumes conditionally independent recognition events; correlated responses require their joint model. This independence must hold under the same conditioning and sampling law used for the probabilities in the product; independence conditional on additional patient characteristics alone does not justify multiplying probabilities averaged over those characteristics. As written, the product model uses the same $r_{ijh}$ across admissible sets $S$, and the union-bound interpretation requires each $q_i$ to remain a valid bound for its corresponding event throughout those sets. These are requirements of that model, not consequences of conditional independence or established biological facts. Test whether increased predicted coverage reduces measured immune escape over a defined horizon.

\paragraph{Additional binding assumptions and resolution criterion.} The displayed minimum over HLA strata is robustness across those strata, not automatically robustness to uncertainty in recognition probabilities. Include a joint confidence or sensitivity set when the latter is claimed. Fix clone weights, or use stratum-specific weights if composition differs by stratum. The cross-reactivity sum is a conservative union-bound constraint only for calibrated compatible event probabilities; finite assay coverage does not establish absence of untested tissue effects.

\paragraph{Alternative formulations.} (i) A personalized formulation optimizes each patient's antigen set using measured clonality and HLA genotype. (ii) A population vaccine seeks shared antigens under a worst-stratum criterion and manufacturing constraints. These are separate design objectives. A dynamic variant incorporates antigen loss and impaired presentation; static coverage is equivalent only when those processes are negligible.

\paragraph{Obstructions and partial progress.} The pancreatic RNA-vaccine study demonstrates immunogenicity in some patients and motivates mechanistic prediction, without establishing general efficacy. Uncertain presentation, immunodominance and shared escape pathways violate naive set-cover assumptions. Antigen availability can vary between metastases, and normal-tissue safety cannot be certified by a finite assay against every tissue state. Vaccine design therefore remains a measured coverage-and-escape problem, distinct from arbitrary combinations of cytotoxic drugs.

\paragraph{Sources (supplied provenance; not all rechecked).} \begin{itemize}
\item Official current archive; verifies the named agenda, not mathematical unsolvability. Unfunded means no team funded through this programme. \url{https://www.cancergrandchallenges.org/challenges/unfunded-challenges}
\item Rojas et al., Personalized RNA neoantigen vaccines stimulate T cells in pancreatic cancer (2023). \url{https://www.nature.com/articles/s41586-023-06063-y} (Primary phase I study; immunogenicity and recurrence associations do not establish universal efficacy.)
\end{itemize}


\section*{Binding catalogue conventions}
This edition revises both supplied catalogues and preserves original problem numbers 1--330. It adds one globally unique label \verb|\label{rank:x}| per problem, where smaller $x$ means greater provisional difficulty. A rank label is an anchor at that problem's section: \verb|\ref{rank:x}| returns the section number, not the literal difficulty number. The visible \texttt{rank:x} text and the index give the difficulty number itself.

The entries are \emph{parameterized mathematical research schemes}, including empirical prediction, causal identification, historical reconstruction, theory and engineering. They are not 330 certified unsolved mathematical conjectures. A rigorous instance fixes the objects and restrictions below; changing those objects changes the question. Established identities, finite calculations, non-identification constructions and solved special cases are not being posed anew as open. Some residual tasks concern informative bounds, biological realization or external validation rather than a new theorem.

Every main definition and problem statement received a mathematical editorial review. Each entry now has an additional binding-conditions paragraph; direct corrections are separately recorded in the accompanying audit. The alternative formulations inherit the applicable conventions, but are different targets unless equivalence is proved. Where an old broad phrase conflicts with an explicit restriction below it, the explicit restriction controls.

\textbf{Literature boundary.} Thirteen selective primary-literature checks cover eighteen entries in this revision; their access levels and scope are recorded explicitly. The supplied bibliography and its earlier source-access descriptions are retained as provenance, not represented as newly verified. Neither this review nor a source's description proves the absence of later solutions. The date of this edition is 13 October 2026; there is no claim of exhaustive current-status certification for all 330 entries.

\subsection*{Binding conventions for a concrete instance}
\paragraph{Objects, quantifiers and non-vacuity.}
An instance declares measurable state, input, output and parameter spaces; a fixed candidate family; admissible initial laws, interventions and environments; a finite horizon or specified asymptotic regime; observation and sampling laws; and performance metrics, positive normalization scales and positive tolerances. These are chosen before fitting or testing, not adjusted after seeing failures. Shared parameters must genuinely be shared across the declared validation domain. Finite dimension alone does not restrict arbitrary real-number encodings: specify a finite library, bounded description length, or regularity/complexity bounds. Exclude constant or no-signal solutions through a fixed informative-output or baseline-improvement requirement where appropriate. A universally quantified result cannot be established by choosing one trivial favorable instance.

\paragraph{Probability laws and data.}
Every expectation or probability refers to the specified stochastic law, including initial uncertainty, process noise, observation noise and sampling as appropriate. A deterministic equation is not a probability model without a declared random input or initial law. ``Observed'' targets in population-level formulas mean population observables $g_j^*$, not noisy estimates asserted to be exact. Let $\widetilde g_j$ be independent validation estimates and suppose
\[
\mathbb{P}\{\lvert\widetilde g_j-g_j^*\rvert\le r_j\text{ for every }j\}\ge1-\alpha.
\]
Then an uncertainty-aware sufficient certification rule for positive target tolerances $\varepsilon_j$ is
\[
\max_j\frac{|\widehat g_j-\widetilde g_j|+r_j}{\varepsilon_j}\le1.
\]
On that simultaneous confidence event it implies every population error is at most its tolerance, by the triangle inequality. Vector/norm versions use corresponding valid simultaneous bounds. Training uncertainty must also be included when the claim averages over repeated training data. Evaluation on finitely many conditions does not establish a continuous-domain supremum without covering and regularity bounds.

\paragraph{Dynamics and well-posedness.}
For ODEs and SDEs give coefficient domains, initial laws and existence/uniqueness conditions in the intended solution sense, or keep the entire admitted solution set. For stochastic jump models specify nonnegative off-diagonal rates, diagonal normalization, rate orientation and non-explosion. For PDEs specify spatial domains, boundary/initial conditions, conservation or source terms, and a weak or classical solution class. Moving-domain concentration equations include dilution/transport terms. A stationary law is used only for a time-homogeneous augmented model under existence assumptions; a stationary approximation to a driven system must be justified. Absorbing extinction states are retained. Survival-conditioned, quasi-stationary and unconditional laws are distinct and must not be substituted for one another.

\paragraph{Identification and causal meaning.}
Two parameter values are observationally equivalent when they induce the same full distribution under all allowed observation/intervention designs, not merely the same likelihood value for one sample. A target is identified only if constant on the equivalence class. Gauge symmetries, label switching and reference conventions are quotiented or fixed. Finite-data compatible sets use a declared confidence, likelihood or discrepancy rule; an empty set signals model rejection. Bayesian credible sets are not automatically confidence sets. A prior or regularizer does not create data identification. Bounds intended to resolve a question must be informative at a predeclared width or decision threshold.

A causal contrast uses well-defined intervention versions and a stated causal model. Random assignment, or explicit consistency, overlap and exchangeability assumptions, is required for causal interpretation from observations. Interference and time-varying confounding are represented when present. Fixing a mediator by intervention differs from conditioning on its observed value. Marginal potential-outcome laws do not determine individual cross-world effects or an arbitrary joint coupling. Ethically admissible designs, rather than unconstrained interventions, define the experimental menu.

\paragraph{Metrics, units, optimization and limits.}
Specify the sampling or reference measure for densities and integrals, the log base, and dimensionless arguments of logarithms. Total variation is $\sup_A|P(A)-Q(A)|$; Wasserstein distances require the declared ground metric and appropriate moment conditions; KL divergence uses its stated absolute-continuity convention. All normed readouts are scaled consistently, and error probabilities use $0<\alpha<1$. A ratio has a positive denominator or an explicit undefined/zero convention. Thermodynamic entropy production is nonnegative in expectation under the relevant hypotheses; individual stochastic realizations need not be. An infimum is not an attained minimum without compactness/coercivity and continuity or another existence argument. Fix the order of large-population, long-time and vanishing-mutation limits. Stability claims declare perturbation space and symmetries.

\subsection*{How the provisional difficulty order was assigned}
Four editorial components are scored from 1 (lower burden) to 5 (higher burden): (a) formal or inferential difficulty, (b) access to evidence and identification burden, (c) mechanistic/computational integration, and (d) independent validation and transfer. Their unweighted total orders the entries. For selected entries in the highest bands, an explicit editorial priority breaks equal totals; all other ties are broken by original problem number. Equal totals are \emph{not} evidence of reliably different true difficulties. Scores are judgments, not measurements, completion-time estimates or consensus rankings.

The object ranked is an informative, nondegenerate resolution of the \emph{scoped} scientific target, with the stated budget and validation domain. A sharp useful impossibility result can qualify; an already known identity, vacuous bound, arbitrary favorable model or lookup table cannot. Rank 330 is merely the least difficult relative assessment in this catalogue, not a claim of an easy biological solution. Broad historical titles do not automatically outrank precisely formulated mechanistic tasks.



\end{document}
